% fig1_dataflow_new.tex
% Figure 1, REDESIGN of fig1_dataflow.tex. Iterating here so the original stays
% untouched until the author adopts this one.
%
% WHY A REDESIGN (defects measured on the built PDF of the original):
%   1. OVERLAPS. "extract" printed on top of the per-visit-vectors box; the three
%      italic legend lines under the graph printed into the sliding-windows box and
%      into its incoming arrow; the per-visit arrow ran down THROUGH the region-
%      vectors box. All three are collisions, not tight spacing.
%   2. IT DID NOT SHOW CHECK2HGI. The four levels were four labelled boxes in a
%      stack -- a legend, not a graph. Nothing in the picture said that the levels
%      are climbed BOTTOM-UP, that a GCN runs inside a level and attention pooling
%      moves between levels, that the visit edges are DIRECTED, that places are
%      linked by a Delaunay graph and regions by adjacency, or that the objective
%      lives at the three level BOUNDARIES. Those are the method.
%   3. SIZE. At \resizebox{0.66\textwidth} the 7pt body rendered at ~6.0pt.
%
% GEOMETRY, and why it is what it is. The paper is at exactly 8 pages, so the one
% hard constraint is RENDERED HEIGHT, not width -- height is what costs a page.
%   original: 397.5 x 178.1 pt natural, drawn at 0.66\textwidth = 340.6pt -> 152.6pt tall.
%   this one:  562.7 x 144.3 pt natural, drawn at \textwidth = 516pt   -> 132.3pt tall.
% So it is 20pt SHORTER than the figure it replaces while carrying far more, and the
% scale factor goes from 0.857 to 0.917, i.e. the 7pt body renders at ~6.4pt instead
% of ~6.0pt. Measured, not estimated: both page sizes come from pdfinfo, and the
% eight-page build was re-run with this file in place and still came out at 8 pages.
% ⚠ THIS REQUIRES main.tex TO CHANGE 0.66\textwidth -> \textwidth on the \resizebox
% line. Left at 0.66 this file renders at ~4.2pt and is unreadable.
%
% READING ORDER: strictly left to right, and no arrow ever runs backwards.
%   input -> the four-level graph panel -> two exported tables -> two windows ->
%   one model -> two predictions.
% The four levels are stacked with the CHECK-IN at the bottom because that is the
% direction of the pass (bottom-up), and because it puts `region` ABOVE `check-in`.
% That ordering is load-bearing: the two exports leave the panel at those two
% heights and stay parallel all the way to the two outputs, so the region thread and
% the check-in thread never cross. (Same reason the slide version of this figure,
% presentation/figures/src/c2h_flow.tex, ordered them that way.)
%
% TWO COLOURS AND ONLY TWO. orange = the check-in / semantic thread; blue = the
% region / spatial thread; teal = the two predictions, which is also the colour of
% the heads in fig2_model.tex, so the two figures agree on what teal means.
%
% FIDELITY -- every claim the drawing makes, and its source:
%   - directed visit edges ................. 04_method.tex; apx_h "Build directed
%     user visit sequence". Drawing them undirected would contradict the paper's
%     own leak argument, so the pictogram has arrowheads and no cycle. The edge
%     WEIGHT decaying with the time gap is real and is NOT drawn: a weight has no
%     honest glyph at this size, and Section IV-B states it in prose.
%   - visit graph is DISCONNECTED ......... an edge joins consecutive visits OF THE
%     SAME USER, so users are separate components. Two components are drawn.
%   - Delaunay at the place level ......... apx_h "A Delaunay triangulation connects
%     nearby places"; 04_method "Nearby places are linked at the place level".
%   - adjacency at the region level ....... apx_h "Region nodes are connected when
%     their polygons intersect".
%   - 15 input features .................... apx_h: 7 category indicators + sin/cos
%     hour + sin/cos weekday + 4 elapsed-time values.
%   - a GCN inside every level ............ apx_h Step 1 (check-in), Step 3 (place),
%     Step 4 (region). The check-in level runs TWO convolution layers, and the
%     drawing says only "GCN": "directed succession, 2 GCNs" boxes at 101.60pt
%     against a 95.4pt column, so printing the count costs either a wrapped
%     descriptor or figure width, and the count is not a claim the paper makes.
%   - attention pooling between levels .... apx_h Step 2 (check-in->place) and
%     Step 4 (place->region); Checkin2POI.py, four heads, one shared seed query.
%   - area-weighted city, TRAINING ONLY ... apx_h Step 5: the city vector "is not an
%     input to the joint forecast model". The lane says "training only" and is set
%     in a lighter grey than the other three, so the level that leads nowhere
%     recedes and the two exported threads carry the eye instead.
%   - three boundary objectives ........... apx_h eq. E.2: L_CP, L_PR, L_R-Omega,
%     each a real-vs-corrupted discrimination. Named in the note under the panel.
%   - 9x64 and 9x64, never concatenated ... apx_h: "The two modalities remain
%     separate; they are not concatenated into one 9x128 matrix." Two strips, and
%     the sentence between them, because two strips alone do not say it.
%
% ONE DELIBERATE OMISSION, flagged so nobody has to rediscover it. The note under
% the panel names the three boundary objectives and not the two auxiliary terms
% (masked-place reconstruction, weight 0.3, and the place-table anchor, 0.1) that
% Section IV-B does mention and apx_h eq. E.2 defines. The note says "one objective
% per level boundary", which is true of the three and does not claim to be the whole
% loss, so this is incomplete rather than wrong. Putting them back costs one line of
% figure height. It is the author's call.
%
% Standalone-compilable (pdflatex fig1_dataflow_new.tex) AND \input-able: the
% wrapper runs only when \FIGINPUT is not already defined, exactly as fig2 does.

\ifdefined\FIGINPUT\else
  \documentclass[border=4pt]{standalone}
  \usepackage{tikz}
  \usepackage{amsmath,amssymb}
  \usetikzlibrary{arrows.meta,positioning,fit,backgrounds,calc}
  \begin{document}
\fi
% In the paper the model box points at Fig. 2 by label; standalone has no labels.
\ifdefined\FIGINPUT \def\figmodelref{\ref{fig:model}} \else \def\figmodelref{2} \fi

% ---------------------------------------------------------------- pictograms
% One pic per level. They are what turns "four levels" from a caption into a
% picture: each says what the EDGES at that level are, which is the only thing
% that distinguishes the four levels from each other.
\tikzset{
  % check-in level: two DIRECTED paths, disconnected (= two users).
  checkins/.pic={
    \begin{scope}[->, line width=0.45pt, >={Latex[length=1.1mm]},
                  shorten >=1.1pt, shorten <=1.1pt, draw=orange!60!black]
      \draw (-0.70,-0.13) -- (-0.44, 0.17);
      \draw (-0.44, 0.17) -- (-0.13, 0.01);
      \draw ( 0.19, 0.15) -- ( 0.45,-0.15);
      \draw ( 0.45,-0.15) -- ( 0.71, 0.09);
    \end{scope}
    \fill[orange!65!black] (-0.70,-0.13) circle (1.15pt) (-0.44,0.17) circle (1.15pt)
          (-0.13,0.01) circle (1.15pt) (0.19,0.15) circle (1.15pt)
          (0.45,-0.15) circle (1.15pt) (0.71,0.09) circle (1.15pt);
  },
  % place level: an undirected triangulation.
  places/.pic={
    \begin{scope}[line width=0.45pt, draw=black!45]
      \draw (-0.62,-0.15) -- (-0.20, 0.18) -- ( 0.22,-0.16) -- ( 0.66, 0.14);
      \draw (-0.62,-0.15) -- ( 0.22,-0.16);
      \draw (-0.20, 0.18) -- ( 0.66, 0.14);
    \end{scope}
    \fill[black!60] (-0.62,-0.15) circle (1.15pt) (-0.20,0.18) circle (1.15pt)
          (0.22,-0.16) circle (1.15pt) (0.66,0.14) circle (1.15pt);
  },
  % region level: polygons that touch, so the edge is a shared border.
  regions/.pic={
    \begin{scope}[line width=0.5pt, draw=blue!55!black, fill=blue!8]
      \draw (-0.62,-0.16) rectangle (-0.22,0.16);
      \draw (-0.14,-0.16) rectangle ( 0.26,0.16);
      \draw ( 0.34,-0.16) rectangle ( 0.74,0.16);
    \end{scope}
    \begin{scope}[line width=0.45pt, draw=blue!55!black]
      \draw (-0.22,0) -- (-0.14,0);
      \draw ( 0.26,0) -- ( 0.34,0);
    \end{scope}
  },
  % city level: one node holding all of them.
  city/.pic={
    \draw[line width=0.5pt, draw=black!45, fill=black!4, rounded corners=1.5pt]
         (-0.52,-0.16) rectangle (0.52,0.16);
    \fill[black!45] (-0.26,0) circle (1.0pt) (0,0) circle (1.0pt) (0.26,0) circle (1.0pt);
  },
}

\begin{tikzpicture}[
    >={Stealth[length=1.9mm]},
    font=\scriptsize,
    % 3.35cm = 95.4pt, and it is a MEASURED width, not a guess: the longest
    % descriptor, "directed succession, GCN", boxes at 90.99pt in \scriptsize\itshape.
    % Anything narrower wraps it, and a wrapped descriptor makes the check-in lane
    % three lines tall, which is what made two lanes touch in the earlier drafts.
    lvl/.style={align=left, text width=3.35cm, inner sep=0pt, font=\scriptsize},
    tag/.style={font=\scriptsize\itshape, text=black!68, align=center, inner sep=1pt},
    spine/.style={->, line width=0.7pt, draw=black!62,
                  >={Stealth[length=2.1mm]}},
    sname/.style={font=\scriptsize\itshape, text=black!68, anchor=west, inner sep=1pt},
    flow/.style={->, line width=0.7pt, draw=black!75},
    oflow/.style={->, line width=0.8pt, draw=orange!65!black},
    bflow/.style={->, line width=0.8pt, draw=blue!55!black},
    ocell/.style={draw=orange!65!black, fill=orange!14, line width=0.45pt,
                  minimum width=1.9mm, minimum height=3.2mm, inner sep=0pt},
    bcell/.style={draw=blue!55!black, fill=blue!12, line width=0.45pt,
                  minimum width=1.9mm, minimum height=3.2mm, inner sep=0pt},
    pred/.style={draw=teal!55!black, fill=teal!10, rounded corners=2.5pt,
                 line width=0.6pt, align=center, inner sep=3pt,
                 minimum width=18.5mm, minimum height=5.5mm},
    model/.style={draw=black!65, dashed, rounded corners=3pt, line width=0.7pt,
                  fill=black!3, align=center, inner sep=3pt,
                  minimum width=24.5mm, minimum height=15mm},
  ]

  % Lane heights, 0.92cm apart: a two-line level descriptor fits between lanes, and
  % four lanes still leave the whole picture shorter than the figure it replaces.
  % Every descriptor must therefore stay at TWO lines -- a third line makes the
  % lanes touch, which is the class of defect this redesign exists to remove.
  \def\yC{1.84}\def\yR{0.92}\def\yP{0.00}\def\yK{-0.92}
  \def\xsp{6.95}\def\xpic{9.45}

  % ------------------------------------------------------------------ the panel
  % Drawn before the two lane tints: both live on the background layer, and within
  % a layer TikZ paints in code order, so a panel written later would cover them.
  \begin{scope}[on background layer]
    \node[draw=black!45, rounded corners=3pt, line width=0.5pt, fill=black!2,
          inner sep=0pt, fit={(3.21,-1.44) (10.51,2.34)}] (panel) {};
    % The two EXPORTED lanes carry a tint, and it is the tint their thread keeps to
    % the end of the picture. Nothing else in the panel is coloured, so the eye
    % finds the two exported levels without reading a word.
    \fill[blue!6,   rounded corners=2pt] (3.31,\yR-0.40) rectangle (10.41,\yR+0.40);
    \fill[orange!8, rounded corners=2pt] (3.31,\yK-0.40) rectangle (10.41,\yK+0.40);
  \end{scope}

  % ===================================================== 1. the raw input
  \foreach \i in {0,1,2,3}{\fill[orange!65!black] (0.30+\i*0.55,\yK) circle (1.2pt);}
  \foreach \i in {0,1,2}{%
    \draw[->, line width=0.5pt, >={Latex[length=1.2mm]}, draw=orange!60!black,
          shorten >=1.8pt, shorten <=1.8pt] (0.30+\i*0.55,\yK) -- (0.85+\i*0.55,\yK);}
  % 2.80cm, and it was measured down to 2.60 and put back: at 2.60 the first label
  % breaks into FOUR lines, which reads as a paragraph rather than a caption.
  \node[tag, text width=2.80cm, anchor=south] at (1.10,\yK+0.16)
       {category, cyclic time,\\elapsed time: 15 values};
  \node[tag, text width=2.80cm, anchor=north] at (1.10,\yK-0.16)
       {one user's visits,\\in time order};
  \draw[oflow] (2.20,\yK) -- (3.26,\yK);

  % ========================================== 2. the four-level graph panel
  % The left column names the level and its EDGES: the edges are the only thing
  % that tells the four levels apart, so they are what the descriptor carries.
  % The two exported levels are named in their thread's colour.
  % The city level is the one that leads NOWHERE: apx_h Step 5 says its vector "is
  % not an input to the joint forecast model". Lighter grey plus the words, so the
  % eye spends its attention on the two lanes that do feed the model.
  \node[lvl] (nC) at (4.99,\yC)
       {\textcolor{black!55}{\textbf{city}}\\\textcolor{black!55}{\textit{one node, training only}}};
  \node[lvl] (nR) at (4.99,\yR)
       {\textcolor{blue!55!black}{\textbf{region}}\\\textit{adjacency graph, GCN}};
  \node[lvl] (nP) at (4.99,\yP) {\textbf{place}\\\textit{Delaunay graph, GCN}};
  \node[lvl] (nK) at (4.99,\yK)
       {\textcolor{orange!65!black}{\textbf{check-in}}\\\textit{directed succession, GCN}};

  \pic at (\xpic,\yC) {city};
  \pic at (\xpic,\yR) {regions};
  \pic at (\xpic,\yP) {places};
  \pic at (\xpic,\yK) {checkins};

  % The bottom-up spine. Each arrow is labelled with WHAT MOVES the data across that
  % boundary; the label sits BESIDE the arrow, never on it, so neither is broken.
  \draw[spine] (\xsp,\yK+0.22) -- (\xsp,\yP-0.22);
  \draw[spine] (\xsp,\yP+0.22) -- (\xsp,\yR-0.22);
  \draw[spine] (\xsp,\yR+0.22) -- (\xsp,\yC-0.22);
  \node[sname] at (\xsp+0.14,\yK+0.46) {attention pooling};
  \node[sname] at (\xsp+0.14,\yP+0.46) {attention pooling};
  \node[sname] at (\xsp+0.14,\yR+0.46) {area-weighted};

  \node[anchor=south, font=\scriptsize\bfseries, text=black!75]
       at ([yshift=0.4mm]panel.north)
       {Check2HGI: a four-level graph, trained without task labels};
  % TWO lines, and the break is forced rather than left to the line breaker. This
  % note is the LOWEST thing in the picture, so every line it takes is height the
  % figure spends, and a third line is what pushed the earlier draft over budget.
  % The per-level graph convolution used to live here and now lives in each level's
  % own descriptor, which is what freed the line.
  \node[anchor=north, tag, text width=6.9cm] at ([yshift=-0.4mm]panel.south)
       {one objective per level boundary
        ($\mathcal{L}_{C\!P}$, $\mathcal{L}_{P\!R}$, $\mathcal{L}_{R\Omega}$)\\
        separates real pairs from corrupted ones};

  % ============================ 3. the two exported tables, as two windows of nine
  \draw[bflow] (10.57,\yR) -- (11.27,\yR);
  \draw[oflow] (10.57,\yK) -- (11.27,\yK);
  \foreach \i in {0,...,8}{\node[bcell] at (11.49+\i*0.22,\yR) {};}
  \foreach \i in {0,...,8}{\node[ocell] at (11.49+\i*0.22,\yK) {};}
  % ⚠ THE WORDING HERE IS A FIDELITY FIX, not a style choice. "one vector per region"
  % reads as a window of nine REGIONS, which is false: both windows hold the SAME nine
  % visits, and what differs is how each visit is represented -- by its own check-in
  % vector, or by the vector of the region it falls in ("the same window with each
  % visit represented by the trained vector of its region node", 04_method.tex). Saying
  % it this way is also what makes the semantic/spatial pair in Fig. 2 legible.
  % Both breaks are forced. Left to the line breaker the top label hyphenated
  % "re-gion's" onto a third line, and the third line is height the figure spends.
  % 3.8cm keeps the left edge clear of the panel (10.51) at 10.67.
  \node[tag, anchor=south, text width=3.8cm] at (12.37,\yR+0.24)
       {$\mathbf{z}_r\!\in\!\mathbb{R}^{64}$ of each visit's region\\a window of $9\times64$};
  \node[tag, anchor=north, text width=3.8cm] at (12.37,\yK-0.24)
       {$\mathbf{x}_i\!\in\!\mathbb{R}^{64}$ of each visit\\a window of $9\times64$};
  % The one fact both the original figure and the appendix's prose call out and that a
  % picture of two strips does NOT carry on its own: they never become one 9x128 matrix.
  \node[tag, text width=2.6cm] at (12.20,\yP) {the two windows\\stay separate};

  % ============================================ 4. one model, two predictions
  \node[model] (mdl) at (15.40,\yP)
       {\textbf{single}\\\textbf{multi-task model}\\one forward pass\\(Fig.~\figmodelref)};
  \draw[bflow] (13.25,\yR) -- ++(0.30,0) |- ([yshift=4mm]mdl.west);
  \draw[oflow] (13.25,\yK) -- ++(0.30,0) |- ([yshift=-4mm]mdl.west);

  \node[pred] (outr) at (18.28, 0.55) {next region};
  \node[pred] (outc) at (18.28,-0.55) {next category};
  \draw[flow] (mdl.east) -- ++(0.42,0) |- (outr.west);
  \draw[flow] (mdl.east) -- ++(0.42,0) |- (outc.west);

\end{tikzpicture}

\ifdefined\FIGINPUT\else
  \end{document}
\fi
